% 1.3 章 · 最小二乘与正交化
% 这一章要把"解一个不相容的方程组"拆成三件互不相同的事：
%   1) 正交化——把问题换成列正交的坐标再做，误差从 kappa^2 u 掉到 kappa u；
%   2) 截断——右端项的噪声把高频分量放大 1/sigma，必须扔掉一部分奇异方向；
%   3) 几何——解不唯一时"哪一个是你要的"，由等残差集与扰动扇面说清楚。
% 断言附表由 tools/build_book.py 从实际运行结果生成，\input 进来，不手抄。

\numlabchapter{最小二乘与正交化}{1.3}

\section{这一章要解决什么}

给定 $A\in\mathbb{R}^{m\times n}$（$m\ge n$）与 $b\in\mathbb{R}^m$，方程 $Ax=b$ 一般**无解**。
退而求其次，我们要的是
\[
  x_{\mathrm{LS}}=\arg\min_{x}\|Ax-b\|_2 .
\]
这件事看起来只是"解一个方程"，实际上要同时回答三个互相独立的问题：

\begin{enumerate}
  \item \textbf{用什么坐标解？} 正规方程 $A^{\mathsf T}Ax=A^{\mathsf T}b$ 在数学上完全正确，
        在浮点里却把条件数**平方**；QR/SVD 用正交变换把误差拉回 $\kappa u$ 量级。
  \item \textbf{什么时候该停下来？} 当 $b$ 里混了噪声，$\sigma_i$ 很小的那些方向会把噪声
        放大成 $1/\sigma_i$；截断 SVD 用"少用几个奇异方向"换取稳定。
  \item \textbf{"解"到底指哪一个？} 秩亏时解集是一张仿射子空间，
        "最小范数"是一条额外约定；$\kappa$ 大到一定程度，"解"本身就没有可表示的有效位。
\end{enumerate}

记号：$u=\tfrac12\varepsilon_{\text{mach}}=2^{-53}\approx1.11\times10^{-16}$ 是单位舍入；
$\kappa=\kappa_2(A)=\sigma_1/\sigma_n$（$\sigma_1\ge\cdots\ge\sigma_n>0$ 是奇异值）；
$\hat{\cdot}$ 表示浮点算出来的量；$\|\cdot\|$ 缺省指 $2$ 范数。

\section{最小二乘的正交性刻画}

\begin{definition}[最小二乘解]
$x_{\mathrm{LS}}$ 称为最小二乘解，若 $\|Ax_{\mathrm{LS}}-b\|_2\le\|Ax-b\|_2$ 对所有
$x\in\mathbb{R}^n$ 成立。残差记为 $r=b-Ax_{\mathrm{LS}}$。
\end{definition}

\begin{theorem}[正规方程]
\label{thm:normal}
$x$ 是最小二乘解，当且仅当 $A^{\mathsf T}(b-Ax)=0$，即 $A^{\mathsf T}Ax=A^{\mathsf T}b$。
\end{theorem}

\begin{proof}
对任意 $z\in\mathbb{R}^n$ 展开目标函数：
\[
  \|A(x+z)-b\|_2^2=\|Ax-b\|_2^2+2z^{\mathsf T}A^{\mathsf T}(Ax-b)+\|Az\|_2^2 .
\]
若 $A^{\mathsf T}(Ax-b)=0$，则右端 $\ge\|Ax-b\|_2^2$，故 $x$ 最小。
反之若 $A^{\mathsf T}(Ax-b)=w\ne0$，取 $z=-\epsilon w$，则
$\|A(x+z)-b\|_2^2=\|Ax-b\|_2^2-2\epsilon\|w\|_2^2+O(\epsilon^2)<\|Ax-b\|_2^2$
（$\epsilon>0$ 充分小），与最小性矛盾。
\end{proof}

\begin{corollary}[残差与列空间正交]
\label{cor:orth}
$A^{\mathsf T}r=0$，即 $r\perp\operatorname{range}(A)$；于是
$\|b\|_2^2=\|Ax_{\mathrm{LS}}\|_2^2+\|r\|_2^2$（勾股定理）。
\end{corollary}

\begin{proof}
由定理 \ref{thm:normal}，$A^{\mathsf T}r=0$ 意味着 $r$ 与 $A$ 的每一列正交。
又 $Ax_{\mathrm{LS}}\in\operatorname{range}(A)$，两向量正交，故勾股定理成立。
\end{proof}

\begin{remark}[正交性是三条路线的共同点]
本章的 A 类断言 \texttt{test\_qr\_least\_squares\_residual\_is\_orthogonal} 与
\texttt{test\_all\_solvers\_satisfy\_normal\_equations} 实测：$\kappa$ 从 $10^{2}$ 涨到 $10^{10}$，
$\|A^{\mathsf T}\hat r\|$ 与后向误差刻度 $n\,u\,\|A\|_F(\|b\|+\|A\|_F\|\hat x\|)$ 的比值始终
$\le0.18$。**正交性与 $\kappa$ 无关**——会随 $\kappa$ 变坏的只有解本身。
\end{remark}

\section{为什么不能直接解正规方程}

\begin{theorem}[条件数的平方]
\label{thm:kappa2}
设 $A$ 列满秩，则 $\kappa_2(A^{\mathsf T}A)=\kappa_2(A)^2$。
\end{theorem}

\begin{proof}
设 $A=U\Sigma V^{\mathsf T}$ 是薄 SVD（$\Sigma=\operatorname{diag}(\sigma_1,\dots,\sigma_n)$）。
则 $A^{\mathsf T}A=V\Sigma^2V^{\mathsf T}$ 是 $A^{\mathsf T}A$ 的特征分解，
特征值为 $\sigma_i^2$。对称正定矩阵的 $\kappa_2$ 是最大最小特征值之比，
故 $\kappa_2(A^{\mathsf T}A)=\sigma_1^2/\sigma_n^2=\kappa_2(A)^2$。
\end{proof}

\begin{theorem}[前向误差的两把尺子]
\label{thm:forward}
设 $\hat x$ 是任一后向稳定算法给出的最小二乘解，$\|A\hat x-b\|$ 与残差同量级，
则
\[
  \frac{\|\hat x-x_{\mathrm{LS}}\|}{\|x_{\mathrm{LS}}\|}
  \;\lesssim\; c\,u\,\kappa_2(A)\quad\text{(正交化方法)},\qquad
  \frac{\|\hat x-x_{\mathrm{LS}}\|}{\|x_{\mathrm{LS}}\|}
  \;\lesssim\; c\,u\,\kappa_2(A)^2\quad\text{(正规方程)} .
\]
\end{theorem}

\begin{proof}[证明思路]
把 $x_{\mathrm{LS}}$ 写成 $A^{+}b$，对 $A$ 与 $b$ 的扰动做一阶展开：
$\delta x=-A^{+}\,\delta A\,x_{\mathrm{LS}}+A^{+}\,\delta b+O(\|\delta\|^2)$，
取范数得 $\|\delta x\|\le\|A^{+}\|\,\|\delta A\|\,\|x_{\mathrm{LS}}\|+\|A^{+}\|\,\|\delta b\|$。
正交化方法作用在 $A$ 上，$\|A^{+}\|=\sigma_n^{-1}$，相对误差是 $u\,\|A\|\cdot\sigma_n^{-1}=u\kappa$；
正规方程作用在 $A^{\mathsf T}A$ 上，其条件数是 $\kappa^2$，而右端 $A^{\mathsf T}b$ 的
浮点误差也被 $\|A\|^2$ 放大，两处相乘给出 $\kappa^2u$。
\end{proof}

实测（A 类断言 \texttt{test\_error\_slopes\_are\_two\_and\_one\_on\_log\_log\_axes}）：
在未饱和窗口内双对数斜率分别是 $1.981$（正规方程）、$0.974$（QR）、$0.988$（SVD），
理论值 $2,1,1$ 全部落在 $\pm0.1$ 内；同一问题上
$\mathrm{err}_{\text{normal}}/\mathrm{err}_{\text{qr}}$ 实测 $1.9\times10^{4}$（$\kappa=10^{4}$）、
$1.6\times10^{7}$（$10^{6}$）、$7.9\times10^{7}$（$10^{8}$）——量级就是 $\kappa$。

\begin{remark}[残差小不等于解准]
定理 \ref{thm:forward} 里的"后向稳定"是**共同前提**：正规方程与 QR 的残差都能小到 $u\|A\|\|x\|$，
但前向误差差 $\kappa$ 倍。本章 \texttt{normal\_vs\_qr} 的主轴同时画误差与残差，
就是为了让这两条曲线在同一个对数坐标系里分道扬镳。
\end{remark}

\section{QR 分解与列主元判秩}

\begin{theorem}[Householder QR]
\label{thm:qr}
任意 $A\in\mathbb{R}^{m\times n}$（$m\ge n$）都可写成 $A=QR$，
其中 $Q\in\mathbb{R}^{m\times m}$ 正交，$R\in\mathbb{R}^{m\times n}$ 上三角。
本库返回**薄**形式：$Q$ 取 $m\times k$、$R$ 取 $k\times n$，$k=\min(m,n)$。
\end{theorem}

\begin{proof}
对 $j=1,\dots,k$，取 $x=A_{j:m,j}$（第一步后是当前子列），
令 $v=x-\alpha e_1$，$\alpha=-\operatorname{sign}(x_1)\|x\|$，
$H_j=I-2vv^{\mathsf T}/\|v\|^2$ 是对称正交矩阵（$H_j^2=I$）且把 $x$ 映到 $\alpha e_1$。
逐列作用 $H_k\cdots H_1A=R$，则 $A=H_1\cdots H_kR=QR$，$Q$ 是正交矩阵之积。
符号取 $-\operatorname{sign}(x_1)$ 是为了让 $v_1=x_1-\alpha=x_1+\operatorname{sign}(x_1)\|x\|$
成为**加法**，从而不丢有效位。
\end{proof}

实测（A 类断言 \texttt{test\_householder\_qr\_reconstruction\_and\_orthogonality}）：
8 种形状（含 $m<n$）上 $\|A-QR\|_F/(n\,u\,\|A\|_F)\le0.94$、
$\max|Q^{\mathsf T}Q-I|/(n\,u)\le0.94$；运算量与闭式 $2n^2(m-n/3)$ **逐位相等**。

\begin{theorem}[列主元 QR 能给的两条界]
\label{thm:rrqr}
设 $\Pi A=QR$ 是列主元 QR（$|r_{11}|\ge|r_{22}|\ge\cdots$），则
\[
  |r_{11}|\;\ge\;\frac{\sigma_1}{\sqrt n},
  \qquad
  \prod_{j=1}^{k}|r_{jj}|\;\le\;\prod_{j=1}^{k}\sigma_j\quad(k=1,\dots,n).
\]
\end{theorem}

\begin{proof}
第一式：$|r_{11}|$ 是 $A$ 的最大列范数，而
$\|A\|_F^2=\sum_i\sigma_i^2\le n\sigma_1^2$ 与
$\max_j\|a_j\|_2\ge\|A\|_F/\sqrt n$ 给出 $|r_{11}|\ge\|A\|_F/\sqrt n\ge\sigma_1/\sqrt n$。
第二式：$\prod_{j\le k}|r_{jj}|=|\det A_k|$（$A_k$ 是前 $k$ 个主元列组成的方阵，
上三角的行列式就是对角元之积），而 $\det$ 的绝对值等于奇异值之积，
截取 $k$ 列只会让奇异值变小：$|\det A_k|=\prod_{j\le k}\sigma_j(A_k)\le\prod_{j\le k}\sigma_j(A)$。
\end{proof}

\begin{remark}[判秩是启发式，不是定理]
常见说法"$|r_{kk}|\approx\sigma_k$"**不成立**：本章实测
$|r_{kk}|/\sigma_k$ 在秩亏矩阵上可取 $0.90,0.99,0.91,2.07$，
Kahan 矩阵更是经典反例（$|r_{nn}|$ 远大于 $\sigma_n$）。
能用的只有定理 \ref{thm:rrqr} 的两条界，加一句工程约定：
$|r_{kk}|\le\max(m,n)\,u\,|r_{11}|\Rightarrow$ 数值秩 $<k$。
判据里必须带容差，秩是"矩阵加容差"的性质。
\end{remark}

\section{奇异值分解}

\begin{theorem}[SVD 的存在性]
\label{thm:svd}
任意 $A\in\mathbb{R}^{m\times n}$ 可写成 $A=U\Sigma V^{\mathsf T}$，
$U\in\mathbb{R}^{m\times m}$、$V\in\mathbb{R}^{n\times n}$ 正交，
$\Sigma$ 非负对角（前 $\min(m,n)$ 个位置）。
\end{theorem}

\begin{proof}
$A^{\mathsf T}A$ 对称半正定，取特征分解 $A^{\mathsf T}A=V\Lambda V^{\mathsf T}$，
$\Lambda=\operatorname{diag}(\lambda_1\ge\cdots\ge\lambda_n\ge0)$，$\sigma_i=\sqrt{\lambda_i}$。
对 $\sigma_i>0$ 定义 $u_i=Av_i/\sigma_i$，则
$u_i^{\mathsf T}u_j=v_i^{\mathsf T}A^{\mathsf T}Av_j/(\sigma_i\sigma_j)
=\lambda_jv_i^{\mathsf T}v_j/(\sigma_i\sigma_j)=\delta_{ij}$：
正奇异值对应的 $u_i$ 已经标准正交，补全成 $\mathbb{R}^m$ 的标准正交基即可，
于是 $AV=U\Sigma$，即 $A=U\Sigma V^{\mathsf T}$。
\end{proof}

\begin{theorem}[Eckart--Young--Mirsky：截断是最优的低秩逼近]
\label{thm:eym}
记 $A_k=\sum_{i\le k}\sigma_iu_iv_i^{\mathsf T}$，则
$\|A-A_k\|_2=\sigma_{k+1}$，且对任意 $\operatorname{rank}(B)\le k$ 有
$\|A-B\|_2\ge\sigma_{k+1}$。
\end{theorem}

\begin{proof}
第一式直接计算：$A-A_k=\sum_{i>k}\sigma_iu_iv_i^{\mathsf T}$，其最大奇异值是 $\sigma_{k+1}$。
第二式用维数论证：设 $\operatorname{rank}(B)\le k$，则 $\ker B$ 的维数 $\ge n-k$，
与 $\operatorname{span}(v_1,\dots,v_{k+1})$ 必有非零交向量 $z$（维数之和 $>n$），
取 $\|z\|=1$ 得
$\|A-B\|_2^2\ge\|(A-B)z\|_2^2=\|Az\|_2^2=\sum_{i\le k+1}\sigma_i^2(v_i^{\mathsf T}z)^2\ge\sigma_{k+1}^2$。
\end{proof}

\begin{theorem}[伪逆与最小范数解]
\label{thm:pinv}
设 $A=U\Sigma V^{\mathsf T}$，$A^{+}=V\Sigma^{+}U^{\mathsf T}$（正奇异值取倒数、零保持零），
则所有最小二乘解中 $x=A^{+}b$ 的范数最小。
\end{theorem}

\begin{proof}
把 $b$ 按 $U$ 的列展开，$\|Ax-b\|_2^2=\sum_i(\sigma_iz_i-c_i)^2+\sum_{i>r}c_i^2$，
其中 $z=V^{\mathsf T}x$、$c=U^{\mathsf T}b$、$r=\operatorname{rank}(A)$：
后一项与 $x$ 无关，前一项在 $z_i=c_i/\sigma_i\ (i\le r)$ 时取零。
自由度只剩 $z_{r+1},\dots,z_n$，它们只会增加 $\|x\|_2=\|z\|_2$，故取零最小。
\end{proof}

实测：四个 Moore--Penrose 条件归一化后最大 $0.39$；
$\|A^{+}-A^{+}_{\text{LAPACK}}\|_{\max}\le5.0\times10^{-16}$；
$\kappa$ 从 $10^{2}$ 到 $10^{14}$ 时 $\sigma_{\min}$ 的相对误差随 $\kappa u$ 线性增长
（$10^{14}$ 时 $8.6\times10^{-4}$）——**大 $\kappa$ 下小奇异值只有相对精度，没有绝对精度**。

\section{两种算法：Jacobi 与 Golub--Kahan}

本库实现了两条独立路线，互为对拍：

\begin{itemize}
  \item \textbf{单边 Jacobi}：反复对 $A$ 的列对 $(p,q)$ 施加右 Givens 旋转，
        使两列正交；旋转角满足 $\cot(2\theta)=(\|a_q\|^2-\|a_p\|^2)/(2a_p^{\mathsf T}a_q)$。
        收敛后用列范数作 $\sigma_i$、归一化列作 $v_i$。
        它几乎不需要证明收敛性之外的技巧，但轮数多、只在 $n$ 小时好用。
  \item \textbf{Golub--Kahan}：先用 Householder 双对角化把 $A$ 化成只有对角与一条次对角线
        非零的 $B$（$A=U_1BV_1^{\mathsf T}$），再对 $B$ 做隐式位移 QR
        （Wilkinson 位移取自尾部 $2\times2$ 块），最后回代累积变换。
        运算量 $O(mn^2)$，是 LAPACK 的做法。
\end{itemize}

\begin{remark}[宽矩阵的双对角形式]
$m<n$ 时 $r\times r$ 的**上**双对角形式不存在：最后一行的尾巴经反射只能变成
$(\beta,0,\dots,0)$，那个 $\beta$ 落在第 $r$ 列之外。本库在 $m<n$ 时转置处理，
返回 \texttt{orientation="lower"} 的**下**双对角 $B$，并显式声明朝向——
这是实现里真踩过的坑（曾导致 $10^{14}$ 倍的重建误差）。
\end{remark}

\section{截断：Picard 条件、L 曲线与偏差原则}

\begin{theorem}[Picard 条件]
\label{thm:picard}
设 $b=\sum_ic_iu_i$，则最小二乘解的系数是 $c_i/\sigma_i$。
若 $b$ 由"精确右端 $b_{\text{exact}}=Ax^{\dagger}$ + 噪声 $\eta$"组成，
$\|\eta\|\approx\delta$，则当 $i$ 大到 $\sigma_i\ll\delta/\|x^{\dagger}\|$ 时，
$c_i/\sigma_i$ 主要由噪声贡献，且随 $i$ 增大而**上升**。
\end{theorem}

\begin{proof}
$b=Ax^{\dagger}+\eta$，而 $Ax^{\dagger}=\sum_i\sigma_i(u_i^{\mathsf T}x^{\dagger})u_i$，
故 $c_i=\sigma_i(u_i^{\mathsf T}x^{\dagger})+u_i^{\mathsf T}\eta$。
当 $\sigma_i|u_i^{\mathsf T}x^{\dagger}|\ll|u_i^{\mathsf T}\eta|\le\|\eta\|$ 时，
$|c_i/\sigma_i|\approx|u_i^{\mathsf T}\eta|/\sigma_i$，随 $\sigma_i$ 下降而上升。
\end{proof}

实测（\texttt{test\_discrepancy\_principle\_bounds\_the\_optimal\_truncation}）：
在 64 组（四种谱形状 $\times$ 四个 $\kappa$ $\times$ 四个噪声水平）配置上，
"偏差原则"给出的截断阶 $k_{\text{dp}}=\#\{\sigma_i>\delta\|b\|/\|x^{\dagger}\|\}$
**从不早于**真实最优：$k_{\text{opt}}\le k_{\text{dp}}+1$ 全部成立，最大超出 $0$。
它是安全的上界，不是最优——L 曲线拐点与 GCV 最小值则可能偏早或偏晚
（实测 cliff 谱上 $k_{\text{opt}}=5$ 而 $k_{\text{cv}}=10$、$k_{\text{corner}}=8$）。

\section{几何：投影、等残差椭圆与噪声扇面}

\begin{theorem}[等残差集是椭球]
\label{thm:ellipse}
设 $x_{\mathrm{LS}}=A^{+}b$、$r_{\min}=b-Ax_{\mathrm{LS}}$，
则 $\{x:\|Ax-b\|_2=\rho\}$（$\rho>\|r_{\min}\|$）是以 $x_{\mathrm{LS}}$ 为中心、
半轴为 $\sqrt{\rho^2-\|r_{\min}\|^2}/\sigma_i$、方向为 $v_i$ 的椭球。
\end{theorem}

\begin{proof}
$\|Ax-b\|_2^2=\|A(x-x_{\mathrm{LS}})\|_2^2+\|r_{\min}\|_2^2$
（交叉项 $2(A(x-x_{\mathrm{LS}}))^{\mathsf T}r_{\min}=0$ 由推论 \ref{cor:orth}）。
令 $y=V^{\mathsf T}(x-x_{\mathrm{LS}})$，则 $\|A(x-x_{\mathrm{LS}})\|_2^2=\sum_i\sigma_i^2y_i^2$，
于是 $\|Ax-b\|_2=\rho$ 等价于 $\sum_i(y_i/(\sqrt{\rho^2-\|r_{\min}\|^2}/\sigma_i))^2=1$，
即一个沿 $v_i$ 的半轴为 $\sqrt{\rho^2-\|r_{\min}\|^2}/\sigma_i$ 的椭球。
\end{proof}

\begin{theorem}[右端扰动在解上的放大]
\label{thm:fan}
设 $b$ 扰动到 $b+\delta b$，$\|\delta b\|=\rho$，
则 $\|\delta x\|_2\le\kappa_2(A)\,\rho/\sigma_1$，且当 $\delta b\parallel u_n$ 时等号成立。
\end{theorem}

\begin{proof}
$\delta x=A^{+}\delta b$，$\|A^{+}\|_2=1/\sigma_n$，故 $\|\delta x\|\le\rho/\sigma_n
=\kappa_2(A)\rho/\sigma_1$；取 $\delta b=\rho u_n$ 得 $\delta x=\rho v_n/\sigma_n$，达到上界。
\end{proof}

实测：$\max\|\delta x\|/(\kappa\delta)=0.281$（三个 $\kappa$ 上完全一致），
最短半轴 $0.28\,\delta$——上界可达但不常达到。
更要紧的是**分辨极限**：解本身自带 $\kappa u$ 的相对舍入，
$\kappa=10^{10}$ 时 $\kappa u\approx1.1\times10^{-6}$ 已经大于默认噪声 $10^{-8}$，
此时扇面量到的是舍入而不是噪声（最短半轴从 $0.28\delta$ 涨到 $8.7\delta$）。
帧里的 \texttt{meta.resolvable} 就是这条判据。

\begin{remark}[病态到一定程度，正规方程不是"算不准"，而是"解不出来"]
$\kappa(A)\gtrsim u^{-1/2}\approx10^{8}$ 之后，$A^{\mathsf T}A$ 的条件数超过 $u^{-1}$：
它的两个特征值在浮点里已经分不开，$A^{\mathsf T}A$ 变成一个**数值奇异**矩阵。
本章的可视化里这不是一句空话：$\kappa=10^{12}$ 时 $3\times2$ 的例子上，
$A^{\mathsf T}A$ 的 LU 回代遇到恰好为 $0$ 的对角元而抛错，
正规方程这条路线**没有有限解**；同一帧里 QR 与 SVD 照常给值，
它们的正交性缺陷 $\|A^{\mathsf T}\hat r\|$ 约为 $10^{-4}\approx\kappa u$。
所以帧里用 \texttt{null}（Python 的 \texttt{None}）而不是 $0$ 或 \texttt{NaN}
表示"没有有限值"，并在 \texttt{strings} 里说明原因：
把"没有数"画成"数很小"或"数很大"都是在编数据。
\end{remark}

\section{它会怎样骗你（本章版）}

\begin{enumerate}
  \item \textbf{把残差小读成解准}：正规方程与 QR 的残差都是 $u\|A\|\|x\|$ 量级，
        前向误差却差 $\kappa$ 倍。实测 $\kappa=10^{8}$ 时 $\|A^{\mathsf T}\hat r\|$ 仍 $\le0.18$ 倍刻度。
  \item \textbf{在饱和段上拟合斜率}：误差到 $O(1)$ 之后幂律就断了，
        把整段塞进拟合会把斜率从 $1.98$ 拉到 $1.52$——拟合窗口必须写明。
  \item \textbf{尺子选错}：同一个 $\|A^{\mathsf T}\hat r\|$，除以 $n\,u\,\|A\|\|b\|$
        在 $\kappa=10^{10}$ 时给出 $5.4\times10^{8}$，除以
        $n\,u\,\|A\|(\|b\|+\|A\|\|\hat x\|)$ 只有 $0.18$：病态时 $\|\hat x\|$ 自己就大。
  \item \textbf{自适应坐标把针画成圆}：等残差椭圆的长短轴之比就是 $\kappa$，
        两轴各自缩放之后 $\kappa=10^{6}$ 的针会变成一个圆，结论当场反转。
  \item \textbf{把 $|r_{kk}|$ 当 $\sigma_k$}：实测比值可以到 $2.07$，Kahan 矩阵更极端。
  \item \textbf{把 L 曲线拐点当最优截断}：实测 cliff 谱上 $k_{\text{corner}}=8$、$k_{\text{opt}}=5$；
        GCV 也会偏（$k_{\text{cv}}=10$）。偏差原则偏保守但方向是安全的。
  \item \textbf{$\kappa u$ 淹没噪声还谈"噪声放大"}：见定理 \ref{thm:fan} 后的实测。
  \item \textbf{把时间轴当时间}：本节各帧的 $t$ 是"扫描进度"或"沿残差滑向 $b$ 的进度"，
        拖到头就夹住，真正要看的是 \texttt{meta} 里的实测值。
\end{enumerate}

% 本文件由 tools/build_book.py 从 dist/reports/assertions.json 生成，请勿手改。
\begin{longtable}{@{}p{0.06\textwidth}p{0.32\textwidth}p{0.22\textwidth}p{0.16\textwidth}@{}}
\caption{章节 1.3 数值断言（生成时间 2026-10-21T12:15:15-0400）}\\
\toprule
类别 & 断言 & 理论值 & 实测值 \\
\midrule
\endfirsthead
\toprule
类别 & 断言 & 理论值 & 实测值 \\
\midrule
\endhead
\bottomrule
\endlastfoot

A & 正规方程斜率 = 2(未饱和窗口) & 2 & 1.98091 \checkmark \\
A & QR 斜率 = 1 & 1 & 0.973565 \checkmark \\
A & 三条曲线的实测斜率 & normal 2 /\allowbreak{} qr 1 /\allowbreak{} svd 1 & normal 1.981,\allowbreak{} qr 0.974,\allowbreak{} svd 0.988(\allowbreak{}拟合窗口上界 $\kappa$ =\allowbreak{} 3.2e+07 /\allowbreak{} 1.0e+12) \checkmark \\
A & err\_\allowbreak{}normal / err\_\allowbreak{}qr(应当 \textasciitilde{} $\kappa$) & 比值 $\in$ [$\kappa$/\allowbreak{}100,\allowbreak{} 100$\kappa$] & $\kappa$=\allowbreak{}1e+04: 1.9e+04；$\kappa$=\allowbreak{}1e+06: 1.6e+07；$\kappa$=\allowbreak{}1e+08: 7.9e+07；实测比值-$\kappa$ 的斜率 0.90 \checkmark \\
A & 与 LAPACK 的相对差(正规方程 /$\kappa$$^2$u，QR 与 SVD /$\kappa$u) & $\le$ c & normal=\allowbreak{}0.394；qr=\allowbreak{}0.657；svd=\allowbreak{}0.396(\allowbreak{}20x8/\allowbreak{}$\kappa$=\allowbreak{}1e+02,\allowbreak{} 30x12/\allowbreak{}$\kappa$=\allowbreak{}1e+04,\allowbreak{} 26x10/\allowbreak{}$\kappa$=\allowbreak{}1e+08) \checkmark \\
A & 最小范数解与 LAPACK 的相对差 & $\|$x\_svd$\|$ =\allowbreak{} $\|$x\_LAPACK$\|$(\allowbreak{}最小范数)，且 $\|$x\_svd$\|$ $\le$ $\|$x\_任意解$\|$ & $\|$x\_svd$\|$=\allowbreak{}3.569390；$\|$x\_lapack$\|$=\allowbreak{}3.569390；$\|$x\_任意解$\|$=\allowbreak{}3.575530；残差 3.96e-15；相对差 4.20e-16 \checkmark \\
A & k\_\allowbreak{}opt $\le$ k\_\allowbreak{}dp + 1 成立的比例 & 全部成立(\allowbreak{}偏差原则偏保守) & 64/\allowbreak{}64；最大超出 0；样例 decay/\allowbreak{}1e4/\allowbreak{}$\delta$=\allowbreak{}1e-4: k\_opt=\allowbreak{}6 $\le$ k\_dp=\allowbreak{}10；decay/\allowbreak{}1e4/\allowbreak{}$\delta$=\allowbreak{}1e-6: k\_opt=\allowbreak{}7 $\le$ k\_dp=\allowbreak{}10；decay/\allowbreak{}1e4/\allowbreak{}$\delta$=\allowbreak{}1e-8: k\_opt=\allowbreak{}10 $\le$ k\_dp=\allowbreak{}10；decay/\allowbreak{}1e4/\allowbreak{}$\delta$=\allowbreak{}1e-10: k\_opt=\allowbreak{}10 $\le$ k\_dp=\allowbreak{}10 \checkmark \\
A & 椭圆上残差的相对起伏(相对 $\rho$) & $\le$ c$\cdot$$\kappa$u + 1e-12 & $\kappa$=\allowbreak{}1e+02:7.4e-13；$\kappa$=\allowbreak{}1e+06:1.6e-10；$\kappa$=\allowbreak{}1e+10:1.4e-06 \checkmark \\
A & max$\|$$\delta$x$\|$/($\kappa$$\cdot$$\delta$)；最短半轴 /$\delta$ & 放大 $\le$ $\kappa$$\delta$(\allowbreak{}比值 $\le$ 1)；$\delta$ $\gg$ $\kappa$u 时最短半轴 $\approx$ $\delta$ & $\kappa$=\allowbreak{}1e+02: max 2.81e-07 =\allowbreak{} 0.281$\cdot$$\kappa$$\delta$,\allowbreak{} min 0.28$\cdot$$\delta$；$\kappa$=\allowbreak{}1e+06: max 2.81e-03 =\allowbreak{} 0.281$\cdot$$\kappa$$\delta$,\allowbreak{} min 0.28$\cdot$$\delta$；$\kappa$=\allowbreak{}1e+08: max 2.81e-01 =\allowbreak{} 0.281$\cdot$$\kappa$$\delta$,\allowbreak{} min 0.25$\cdot$$\delta$；$\kappa$=\allowbreak{}1e10(\allowbreak{}$\kappa$u=\allowbreak{}1.1e-06 $\gg$ $\delta$=\allowbreak{}1e-08)最短半轴 8.66$\cdot$$\delta$ \checkmark \\
B & $\|$A$^{\mathsf T}$r\_\allowbreak{}hat$\|$ /(n$\cdot$u$\cdot$$\|$A$\|$\_\allowbreak{}F($\|$b$\|$+$\|$A$\|$\_\allowbreak{}F$\|$x\_\allowbreak{}hat$\|$)) & $\le$ c(\allowbreak{}三条路线、\allowbreak{}$\kappa$ 从 1e2 到 1e10 都与 $\kappa$ 无关) & 0.180 \checkmark \\
B & $\|$x\_\allowbreak{}k$\|$ 非减(decay) & nondecreasing & 违反 0 处 \checkmark \\
B & 残差非增(decay) & nonincreasing & 违反 0 处 \checkmark \\
B & $\|$x\_\allowbreak{}k$\|$ 非减(cliff) & nondecreasing & 违反 0 处 \checkmark \\
B & 残差非增(cliff) & nonincreasing & 违反 0 处 \checkmark \\
B & $\|$x\_\allowbreak{}k$\|$ 非减(plateau) & nondecreasing & 违反 0 处 \checkmark \\
B & 残差非增(plateau) & nonincreasing & 违反 0 处 \checkmark \\
B & 通过单调性检查的谱形状数 & 3 & 3 \checkmark \\
B & 四种谱的 L 曲线拐点 & 拐点由 log-log 离散曲率最大处给出 & decay:k=\allowbreak{}5；cliff:k=\allowbreak{}8；plateau:k=\allowbreak{}4；two\_clusters:k=\allowbreak{}2 \checkmark \\
B & max|P$^2$$-$P|、max|P$^{\mathsf T}$$-$P|、max$\|$Pv$\|$/$\|$v$\|$ & $\le$ c$\cdot$$\kappa$u(\allowbreak{}$\kappa$=\allowbreak{}1e6 时实测 1.1e-11)、\allowbreak{}0、\allowbreak{}$\le$ 1 & 1.06e-11、\allowbreak{}1.49e-11、\allowbreak{}0.756856 \checkmark \\
B & 端点误差 & t=\allowbreak{}0: (\allowbreak{}p\_t,\allowbreak{}r\_t)=\allowbreak{}(\allowbreak{}p,\allowbreak{}r\_LS)；t=\allowbreak{}1: $\|$r\_t$\|$=\allowbreak{}0、\allowbreak{}p\_t=\allowbreak{}b；任意 t: $\|$p\_t+r\_t$-$b$\|$ $\le$ u$\|$b$\|$ & t=\allowbreak{}0 $\|$r\_t$-$(\allowbreak{}b$-$p\_LS)$\|$=\allowbreak{}0.000e+00；t=\allowbreak{}1 $\|$r\_t$\|$=\allowbreak{}0.000e+00、\allowbreak{}$\|$p\_t$-$b$\|$=\allowbreak{}0.000e+00；t=\allowbreak{}0.5 合成误差 1.241e-16 \checkmark \\
B & 实测 $\kappa$ / 目标 $\kappa$ $-$ 1；四种谱 & $\le$ 1e-6(\allowbreak{}$\sigma$\_min 只精确到 \textasciitilde{}$\kappa$u) & 4.56e-08；decay:$\sigma$\_min=\allowbreak{}1.0e-08；cliff:$\sigma$\_min=\allowbreak{}1.0e-08；plateau:$\sigma$\_min=\allowbreak{}1.0e-08；two\_clusters:$\sigma$\_min=\allowbreak{}1.0e-08 \checkmark \\
B & 通过自检的 (viz, t) 组合数 & 4 个可视化 $\times$ 各自的时间采样点 & 12；projection\_3d:15数组/\allowbreak{}18meta；normal\_vs\_qr:20数组/\allowbreak{}32meta；svd\_truncation:17数组/\allowbreak{}22meta；uncertainty\_fan:17数组/\allowbreak{}25meta \checkmark \\
B & 实测 $\kappa$ 与网格的相对差($\kappa$ 本身只能精确到 \textasciitilde{}$\kappa$u) & $\le$ 1.00000e-06 & $1.551\times 10^{-8}$ \checkmark \\
B & 曲线有限性 /(地板 u) 与 $\kappa$ 一致性 & 全部有限且 $\ge$ u =\allowbreak{} 1.1e-16 & $\kappa$ 相对差 1.55e-08；11 条曲线全部有限；地板 1.110e-16 \checkmark \\
B & GCV 残差随 $\lambda$ 非减 & nondecreasing & 违反 0 处 \checkmark \\
B & GCV 解范数随 $\lambda$ 非增 & nonincreasing & 违反 0 处 \checkmark \\
B & CV 曲线最小值位置与曲线条数 & 最小值在网格内部(\allowbreak{}不是端点) & $\lambda$ $\in$ [1.0e-09,\allowbreak{} 1.0e+01]，k\_cv =\allowbreak{} 8，最小 CV =\allowbreak{} 2.948e-14(\allowbreak{}真最优 k\_opt =\allowbreak{} 5，L 曲线拐点 =\allowbreak{} 5) \checkmark \\
B & 角点帧数(build+pack+有限性) & 4 个 viz $\times$ 参数角点 $\times$ 时间端点 & 94 帧全部通过；分章：normal\_vs\_qr 22、\allowbreak{}projection\_3d 23、\allowbreak{}svd\_truncation 23、\allowbreak{}uncertainty\_fan 26 \checkmark \\
B & $\kappa$=1e12 角点：帧可编码 + 奇异被标注 / $\kappa$=1e3 上解仍正确 & 奇异 $\to$ meta 用 None/\allowbreak{}True + strings 说明；健康点 $\to$ 解与 pinv 一致(\allowbreak{}< 1e-10) & 3 种 method $\times$ 3 个 t 全部编码成功；selected=\allowbreak{}normal 时 3 个 meta 均为 None；$\kappa$=\allowbreak{}1e3 的 QR 投影相对差 1.1e-14 \checkmark \\
A & 重建误差 /(n$\cdot$u$\cdot$$\|$A$\|$\_\allowbreak{}F) & $\le$ 1 & 0.940(\allowbreak{}各形状 2x2:4.4e-16,\allowbreak{} 8x5:2.2e-15,\allowbreak{} 5x8:1.3e-15,\allowbreak{} 12x12:1.8e-15,\allowbreak{} 20x4:1.9e-15,\allowbreak{} 30x30:6.2e-15,\allowbreak{} 40x12:3.1e-15,\allowbreak{} 6x40:2.7e-15) \checkmark \\
B & 正交性缺陷 /(n$\cdot$u) & $\le$ 10 & 2 \checkmark \\
A & $\|$Hx$-$$\alpha$e$_1$$\|$/$\|$x$\|$(稳定取号 vs 朴素取号) & 稳定 \textasciitilde{}u；朴素 \textasciitilde{}$\|$x\_perp$\|$/\allowbreak{}$\|$x$\|$ & $\varepsilon$=\allowbreak{}1e-02: 稳定 4.4e-16 /\allowbreak{} 朴素 1.0e-02(\allowbreak{}横向 1.1e-02)；$\varepsilon$=\allowbreak{}1e-05: 稳定 2.3e-21 /\allowbreak{} 朴素 1.0e-05(\allowbreak{}横向 1.1e-05)；$\varepsilon$=\allowbreak{}1e-08: 稳定 6.6e-26 /\allowbreak{} 朴素 1.0e-08(\allowbreak{}横向 1.1e-08)；$\varepsilon$=\allowbreak{}1e-11: 稳定 6.6e-35 /\allowbreak{} 朴素 1.0e-11(\allowbreak{}横向 1.1e-11) \checkmark \\
B & 正交性 Q$^{\mathsf T}$Q = I & $\le$ 2.66454e-15 & $3.331\times 10^{-16}$ \checkmark \\
A & 重建误差 /(n$\cdot$u$\cdot$$\|$A$\|$\_\allowbreak{}F)，旋转个数 & $\le$ 1；旋转数 =\allowbreak{} $\Sigma$\_\{col<k\}(\allowbreak{}m$-$1$-$col) & 最大 0.837；个数 2x2:1,\allowbreak{} 8x5:25,\allowbreak{} 5x8:10,\allowbreak{} 12x12:66,\allowbreak{} 20x4:70,\allowbreak{} 30x30:435,\allowbreak{} 40x12:402,\allowbreak{} 6x40:15(\allowbreak{}全部等于闭式计数) \checkmark \\
A & $\|$A$^{\mathsf T}$r$\|$ /(n$\cdot$u$\cdot$$\|$A$\|$\_\allowbreak{}F$\cdot$($\|$b$\|$+$\|$A$\|$\_\allowbreak{}F$\|$x\_\allowbreak{}hat$\|$)) & $\le$ c(\allowbreak{}与 $\kappa$ 无关) & 0.180(\allowbreak{}$\kappa$=\allowbreak{}1e+02:0.042,\allowbreak{} $\kappa$=\allowbreak{}1e+06:0.038,\allowbreak{} $\kappa$=\allowbreak{}1e+10:0.180)；与 LAPACK 解的最大相对差 2.87e-08 \checkmark \\
A & 秩亏 $\to$ 数值秩 3；近秩亏 $\to$ 数值秩 4 & 判秩用 |r\_nn| $\le$ max(\allowbreak{}m,\allowbreak{}n)u|r$_1$$_1$|(\allowbreak{}启发式)；可证明的是 |r$_1$$_1$| $\ge$ $\sigma$$_1$/\allowbreak{}$\surd$n 与体积界(\allowbreak{}含舍入松弛) & 秩亏 rank=\allowbreak{}3(\allowbreak{}|r$_4$$_4$|=\allowbreak{}7.35e-16 $\le$ tol=\allowbreak{}8.57e-15)，近秩亏 rank=\allowbreak{}4；|r\_kk|/\allowbreak{}$\sigma$\_k=\allowbreak{}[0.9  0.99 0.91 2.07]；|r\_kk|/\allowbreak{}$\sigma$\_k=\allowbreak{}[0.9  0.99 0.91 1.22] \checkmark \\
A & n\_\allowbreak{}ops 与 qr\_\allowbreak{}flops 的最大差 & 0(\allowbreak{}逐位相等) & 0；实测 4x4:160,\allowbreak{} 10x6:860,\allowbreak{} 30x30:47740,\allowbreak{} 40x12:12736 \checkmark \\
A & max\_\allowbreak{}i |$\Delta$$\sigma$\_\allowbreak{}i|/$\sigma$\_\allowbreak{}i $\div$ (n$\cdot$u$\cdot$$\sigma$\_\allowbreak{}max/$\sigma$\_\allowbreak{}i) & $\le$ 1 & jacobi=\allowbreak{}0.342；kogbetliantz=\allowbreak{}0.361；golub\_kahan=\allowbreak{}0.271(\allowbreak{}形状 2x2,\allowbreak{} 8x5,\allowbreak{} 5x8,\allowbreak{} 12x12,\allowbreak{} 20x4,\allowbreak{} 30x30,\allowbreak{} 40x12,\allowbreak{} 6x40) \checkmark \\
A & 重建误差 /(n$\cdot$u$\cdot$$\|$A$\|$\_\allowbreak{}F) 与两翼正交性缺陷 /(n$\cdot$u) & 两者都 $\le$ c & 重建比 4.585，正交性比 5.000 \checkmark \\
A & 奇异向量残差 /(u$\cdot$$\|$A$\|$\_\allowbreak{}F) & $\le$ c & 42.08 \checkmark \\
A & 四条定义式的最大缺陷 /(n$\cdot$u) & $\le$ c & 0.392；9x5:3.9e-16,\allowbreak{} 9x5:1.9e-16,\allowbreak{} 14x4:5.6e-16；与 LAPACK pinv 的元素最大差 5.0e-16 \checkmark \\
A & 秩亏矩阵的数值秩；$\kappa$$_2$ 与 LAPACK 的相对差 & 秩 =\allowbreak{} 4；$\kappa$$_2$ 相对差 $\le$ c$\cdot$$\kappa$$\cdot$u & 秩 \{'jacobi': 4,\allowbreak{} 'kogbetliantz': 4,\allowbreak{} 'golub\_kahan': 4\}；$\kappa$$_2$ 实测 1.000000001e+08 vs LAPACK 1.000000001e+08，相对差 6.36e-10 \checkmark \\
A & 双对角化重建误差 /(n$\cdot$u$\cdot$$\|$A$\|$\_\allowbreak{}F) & $\le$ 1；非对角外元素 =\allowbreak{} 0 & 0.803 \checkmark \\
A & max|$\sigma$ $-$ $\sigma$\_\allowbreak{}解析| / (n$\cdot$u$\cdot$$\sigma$\_\allowbreak{}max) & $\le$ 1 & 0.750(\allowbreak{}[[1,\allowbreak{}1],\allowbreak{}[0,\allowbreak{}1]],\allowbreak{} diag(\allowbreak{}3,\allowbreak{}1,\allowbreak{}0.5),\allowbreak{} 4$\cdot$Q(\allowbreak{}4x3)) \checkmark \\
A & 小奇异值相对误差 /(n$\cdot$u$\cdot$$\kappa$) & $\le$ c(\allowbreak{}相对误差 \textasciitilde{} $\kappa$u，不是 (\allowbreak{}$\kappa$u)$^2$) & 最大比 0.040；$\kappa$=\allowbreak{}1e+02: 3.5e-15/\allowbreak{}1.7e-15/\allowbreak{}1.1e-15；$\kappa$=\allowbreak{}1e+06: 9.9e-13/\allowbreak{}2.1e-12/\allowbreak{}1.8e-12；$\kappa$=\allowbreak{}1e+10: 1.3e-07/\allowbreak{}1.9e-08/\allowbreak{}1.1e-08；$\kappa$=\allowbreak{}1e+14: 8.6e-04/\allowbreak{}2.3e-04/\allowbreak{}2.2e-04 \checkmark \\
B & 结构检查通过的形状数 & 8 & 8 \checkmark \\
B & |r\_\allowbreak{}kk| 非增 (12x6) & nonincreasing & 违反 0 处 \checkmark \\
B & 剩余列范数非增 (12x6) & nonincreasing & 违反 0 处 \checkmark \\
B & |r\_\allowbreak{}kk| 非增 (20x8) & nonincreasing & 违反 0 处 \checkmark \\
B & 剩余列范数非增 (20x8) & nonincreasing & 违反 0 处 \checkmark \\
B & |r\_\allowbreak{}kk| 非增 (30x5) & nonincreasing & 违反 0 处 \checkmark \\
B & 剩余列范数非增 (30x5) & nonincreasing & 违反 0 处 \checkmark \\
B & offdiag 范数非增 (8x5) & nonincreasing & 违反 0 处 \checkmark \\
B & offdiag 范数非增 (12x8) & nonincreasing & 违反 0 处 \checkmark \\
B & offdiag 范数非增 (20x6) & nonincreasing & 违反 0 处 \checkmark \\
B & 三种形状扫描 4 次后的 offdiag 范数 & 单调不增 & 起点 1.45e+00,\allowbreak{} 3.19e+00,\allowbreak{} 1.78e+00 \checkmark \\
B & 排序/符号规范化后的奇异向量残差 /(u$\cdot$$\|$A$\|$\_\allowbreak{}F) & $\sigma$ 降序非负；残差 $\le$ c & 21.23 \checkmark \\
B & Givens 矩阵正交 & $\le$ 3.55271e-15 & $2.665\times 10^{-17}$ \checkmark \\
B & Givens 矩阵正交 & $\le$ 3.55271e-15 & 0 \checkmark \\
B & Givens 矩阵正交 & $\le$ 3.55271e-15 & 0 \checkmark \\
B & Givens 矩阵正交 & $\le$ 3.55271e-15 & 0 \checkmark \\
B & Givens 矩阵正交 & $\le$ 3.55271e-15 & 0 \checkmark \\
B & |c$^2$+s$^2$$-$1| 与 |det G $-$ 1| 的最大值 & $\le$ u & 0.00e+00 /\allowbreak{} 0.00e+00 \checkmark \\
B & 回代最大相对误差 & $\le$ n$\cdot$u$\cdot$$\kappa$(\allowbreak{}R) & 1.86e-16 \checkmark \\
B & k=n 的截断解与伪逆解的逐位差 & 0 & 4.44e-16 \checkmark \\
B & Householder 矩阵正交 & $\le$ 6.21725e-15 & $5.551\times 10^{-16}$ \checkmark \\
B & Householder 矩阵正交 & $\le$ 1.33227e-14 & $2.220\times 10^{-16}$ \checkmark \\
B & 列范数最大绝对差 & $\le$ u$\cdot$$\|$a\_j$\|$ & 8.88e-16 \checkmark \\
\end{longtable}

